% Источники: exam/materials/преподаватель/2026/Моделирование_случайных_величин,_векторов_.pdf; exam/materials/прошлогодние ответы/Modelirovanie_1_-3.pdf, вопрос 19; exam/materials/прошлогодние ответы/Ответы_Моделирование.pdf, вопрос 21.
\section{Общие методы моделирования случайных процессов.}

\textbf{Случайный процесс} $X(t)$ — параметрическое семейство случайных величин, где параметром является время $t$. Характеризуется:
\begin{itemize}
  \item математическим ожиданием $m_X(t) = M[X(t)]$;
  \item корреляционной функцией $K_X(t_1,t_2) = M[(X(t_1)-m_X(t_1))(X(t_2)-m_X(t_2))]$.
\end{itemize}

\subsection*{Дискретизация процесса}

Так как случайные процессы представляют собой аналоговую функцию (случайную функцию), то для моделирования случайного процесса на ЦВМ необходима дискретизация процесса. Дискретные значения, полученные с некоторым шагом во времени, можно считать случайным вектором.

\subsection*{Метод условных распределений}

Процесс генерируется последовательно:
\[
  x(t_{k+1}) \sim F\bigl(x_{k+1}\mid x(t_k) = x_k\bigr).
\]
Применим, когда условное распределение известно (марковские процессы, авторегрессии).

\subsection*{Рекуррентный метод}

В преподавательском материале для моделирования стационарных случайных процессов используются рекуррентные алгоритмы. Для процесса с показательным распределением записана схема:
\[
  \xi(t_n) = \rho_n\,\xi(t_{n-1}) + \sqrt{1 - \rho_n^2}\cdot x[n],
\]
где
\[
  \rho_n = e^{-\omega(t_n-t_{n-1})}, \qquad \omega=\frac{\pi}{2\Delta t},
\]
а $x[n]$ — последовательность независимых нормальных случайных величин с параметрами $(0,1)$.

Для моделирования случайных процессов не существует универсального подхода: для генерации процесса с заданным законом распределения подбирают рекуррентную функцию. Если процесс является марковским, применяют методы моделирования марковских процессов.

\clearpage
